\documentclass[11pt]{article}
\usepackage{amsmath,amssymb}
\usepackage[margin=1in]{geometry}
\usepackage{booktabs}
\usepackage{hyperref}
\usepackage{xcolor}

\newcommand{\code}[1]{\texttt{#1}}

\title{AGN X-ray Heating: Derivation of the Emissivity-to-Heating-Rate Pipeline}
\author{Fatma Yousef}
\date{\today}

\begin{document}
\maketitle

\noindent This note derives, from the AGN spectral model down to the $T_K$
evolution equation, every equation implemented in
\code{ComputeTs.c}/\code{XRayHeatingFunctions.c} for the AGN X-ray channel,
in parallel with the pre-existing stellar (GAL) pathway. Every symbol is
cross-referenced to its code variable so this can be checked line-by-line
against the source. The specific question under discussion — why the soft
and hard bands need distinct integration limits in the
\emph{luminosity-conversion} step but not, for the same reason, in
\code{integrate\_over\_nu} — is addressed explicitly in Section~\ref{sec:why-limits}.

\tableofcontents

\section{The AGN spectral model}

Each band (soft, hard) is modelled as an independent power law in
specific luminosity, pivoted at a threshold frequency $\nu_{\rm th}$
(\code{NuXrayThreshold}):
\begin{equation}
  L_b(\nu) = L_{0,b}\left(\frac{\nu}{\nu_{\rm th}}\right)^{-\alpha_b},
  \qquad b \in \{\mathrm{soft, hard}\},
  \label{eq:sed}
\end{equation}
where $\alpha_b$ is \code{SpecIndexXrayAGNSoft} or \code{SpecIndexXrayAGNHard},
and $L_{0,b}$ is the (as yet unknown) amplitude at the pivot. The two bands
are defined over disjoint frequency ranges, split at the break frequency
$\nu_{\rm break}$ (\code{NuXraySoftCut}):
\begin{align}
  \text{soft:} \quad & \nu \in [\nu_{\rm th},\, \nu_{\rm break}], \\
  \text{hard:} \quad & \nu \in [\nu_{\rm break},\, \nu_{\rm max}], \qquad \nu_{\rm max} = \code{NuXrayMax}.
\end{align}

Each band's amplitude $L_{0,b}$ is fixed independently, by requiring that
integrating Eq.~\eqref{eq:sed} over that band's own frequency range
reproduces a physically known, calibrated band luminosity $L_{{\rm band},b}$
(the accretion-derived $L_X$ from the black-hole growth history, computed in
\code{blackhole\_feedback.c} — \code{quasar\_lx}, obscuration-corrected):
\begin{equation}
  \int_{\nu_{{\rm lo},b}}^{\nu_{{\rm hi},b}} L_b(\nu)\,d\nu = L_{{\rm band},b}.
  \label{eq:calibration}
\end{equation}

\section{Step 1: the luminosity conversion factor}
\label{sec:conversion}

Define the \emph{luminosity conversion factor} as the ratio of the pivot
amplitude to the calibrated band luminosity,
\begin{equation}
  C_b \equiv \frac{L_{0,b}}{L_{{\rm band},b}},
\end{equation}
which is exactly \code{Luminosity\_converstion\_factor\_AGN\_soft} /
\code{\_hard} in the code (up to the unit conversions in
Section~\ref{sec:units}). Solving Eq.~\eqref{eq:calibration} for $L_{0,b}$
gives $C_b$.

\subsection{Case $\alpha_b \neq 1$}

\begin{align}
  L_{{\rm band},b} &= L_{0,b}\,\nu_{\rm th}^{\alpha_b}
     \int_{\nu_{{\rm lo},b}}^{\nu_{{\rm hi},b}} \nu^{-\alpha_b}\,d\nu
   = L_{0,b}\,\nu_{\rm th}^{\alpha_b}\,
     \frac{\nu_{{\rm hi},b}^{1-\alpha_b} - \nu_{{\rm lo},b}^{1-\alpha_b}}{1-\alpha_b}, \\[4pt]
  \Rightarrow\quad
  C_b &= \frac{(1-\alpha_b)\,\nu_{\rm th}^{-\alpha_b}}
              {\nu_{{\rm hi},b}^{1-\alpha_b} - \nu_{{\rm lo},b}^{1-\alpha_b}}.
  \label{eq:C-generic}
\end{align}

\noindent \textbf{Code check} (soft band, \code{ComputeTs.c:1027-1037}):
\begin{verbatim}
Luminosity_converstion_factor_AGN_soft =
    pow(nu_hi, 1 - alpha) - pow(nu_lo, 1 - alpha);      // (nu_hi^(1-a) - nu_lo^(1-a))
Luminosity_converstion_factor_AGN_soft = 1.0 / ...;     // 1 / (...)
Luminosity_converstion_factor_AGN_soft *=
    pow(nu_th, -alpha) * (1 - alpha);                   // x nu_th^-a (1-a)
\end{verbatim}
which is exactly Eq.~\eqref{eq:C-generic}, per band, with
$(\nu_{{\rm lo},b},\nu_{{\rm hi},b}) = (\nu_{\rm th},\nu_{\rm break})$ for soft
and $(\nu_{\rm break},\nu_{\rm max})$ for hard.

\subsection{Case $\alpha_b = 1$ (the singular case)}

When $\alpha_b \to 1$, $1-\alpha_b \to 0$ and Eq.~\eqref{eq:C-generic} is
$0/0$: the power-law primitive $\nu^{1-\alpha}/(1-\alpha)$ degenerates,
because $\int \nu^{-1}d\nu = \ln\nu$ instead. Re-deriving directly for
$\alpha_b=1$:
\begin{align}
  L_{{\rm band},b} &= L_{0,b}\,\nu_{\rm th}
     \int_{\nu_{{\rm lo},b}}^{\nu_{{\rm hi},b}} \nu^{-1}\,d\nu
   = L_{0,b}\,\nu_{\rm th}\,\ln\!\left(\frac{\nu_{{\rm hi},b}}{\nu_{{\rm lo},b}}\right), \\[4pt]
  \Rightarrow\quad
  C_b &= \frac{1}{\nu_{\rm th}\,\ln(\nu_{{\rm hi},b}/\nu_{{\rm lo},b})}.
  \label{eq:C-unity}
\end{align}

\noindent \textbf{Code check} (\code{ComputeTs.c:1021-1026}):
\begin{verbatim}
if (fabs(SpecIndexXrayAGNSoft - 1.0) < REL_TOL) {
    Luminosity_converstion_factor_AGN_soft =
        (NuXrayThreshold*NU_over_EV) * log(NuXraySoftCut / NuXrayThreshold);
    Luminosity_converstion_factor_AGN_soft = 1.0 / ...;
}
\end{verbatim}
which is exactly Eq.~\eqref{eq:C-unity}. This confirms the branch is not an
arbitrary numerical safeguard: $\alpha_b=1$ is a genuine mathematical
singularity of the power-law primitive (a log-divergence of the naive
formula), not just a coding edge case, so it needs its own closed-form
solution rather than a tolerance-based patch of Eq.~\eqref{eq:C-generic}.
The same structure (generic $\alpha\neq1$ branch + $\alpha=1$ log branch)
already exists for GAL and, under \code{USE\_MINI\_HALOS}, for PopIII —
the AGN soft/hard code simply repeats it per band.

\subsection{Unit conversion}
\label{sec:units}

$C_b$ as derived above is dimensionless (it is a pure ratio of
frequency-integrated shape to the same shape's own normalization). The code
additionally divides by \code{PLANCK} (Section \code{ComputeTs.c:1042,1066}):
this converts a specific-luminosity amplitude into a specific
\emph{photon}-luminosity amplitude ($L_\nu/h\nu_{\rm th}$-type normalization),
consistent with everything downstream (\code{integrate\_over\_nu}'s
integrand, Section~\ref{sec:shape}) being expressed per unit photon energy.
Unlike the GAL/PopIII conversion factors, no \code{SEC\_PER\_YEAR} factor is
applied: GAL/III start from a star-formation-rate calibration
($M_\odot\,{\rm yr}^{-1}$) and need that factor to convert to a per-second
rate; the AGN band luminosities are already luminosities
(${\rm erg\,s^{-1}}$), so no such conversion is needed.

\section{Step 2: the redshift-dependent prefactor}

The pivot amplitude, once fixed by $C_b$, still needs to be turned into the
proper prefactor multiplying the frequency-shape integral at each emission
redshift $z'$ (\code{zp}). This is \code{const\_zp\_prefactor\_AGN\_soft}/
\code{\_hard}:
\begin{equation}
  \mathcal{A}_b(z') = \frac{C_b}{\nu_{\rm th}}\, c\,(1+z')^{\alpha_b + 3}.
  \label{eq:prefactor}
\end{equation}
\textbf{Code check} (\code{ComputeTs.c:1117-1131}):
\begin{verbatim}
const_zp_prefactor_AGN_soft = Luminosity_converstion_factor_AGN_soft
    / (NuXrayThreshold * NU_over_EV) * SPEED_OF_LIGHT
    * pow(1 + zp, SpecIndexXrayAGNSoft + 3);
\end{verbatim}
which is Eq.~\eqref{eq:prefactor} directly. This functional form —
$(1+z')^{\alpha+3}$ — is inherited unmodified from the pre-existing GAL
prefactor (itself the standard 21cmFAST/\textsc{Ts.c} form of Pritchard \&
Furlanetto 2007 / Mesinger, Furlanetto \& Cen 2011): the $+3$ comes from
proper number-density dilution $\propto(1+z')^3$, and the $\alpha_b$ comes
from redshifting a $\nu^{-\alpha_b}$ spectrum (each photon's frequency, and
hence its place on the power law, redshifts along with $z'$). The AGN
channel does not re-derive this — it reuses the identical cosmological
scaling, per band, with that band's own $\alpha_b$ and $C_b$.

Both prefactors are checked for finiteness immediately
(\code{ComputeTs.c:1121,1132}) and the run aborts rather than silently
propagating a bad value — this guards against a zero-width band
($\nu_{{\rm hi},b}=\nu_{{\rm lo},b}$, e.g. \code{NuXraySoftCut==NuXrayThreshold}),
which would make $C_b$ divide by zero.

\section{Step 3: the frequency-shape integral (\code{integrate\_over\_nu})}
\label{sec:shape}

This is a \emph{different} integral, doing a \emph{different} job. Where
$C_b$ (Section~\ref{sec:conversion}) fixes one number — the overall
amplitude of the band, so that the model matches a known calibration
luminosity — \code{integrate\_over\_nu} computes, at every radial shell
$R$ (i.e. every look-back shell $z''$, \code{zpp}) and every ionization
state $x_e$, the deposition-efficiency-weighted photon-number shape kernel:
\begin{equation}
  \mathcal{K}_b(z',z'',x_e) = \int_{\nu_{{\rm lo}}(z',z'')}^{\nu_{{\rm hi}}(z',z'')}
    f_{\rm dep}(\nu,x_e)\left(\frac{\nu}{\nu_{\rm th}}\right)^{-\alpha_b-1} d\nu,
  \label{eq:shape}
\end{equation}
where $f_{\rm dep}$ is the frequency- and ionization-dependent
deposition/ionization/Ly$\alpha$-production efficiency (\code{interp\_fheat},
etc. inside \code{integrand\_in\_nu\_heat\_integral} and its ionization/Ly$\alpha$
siblings) and the extra power of $\nu^{-1}$ relative to Eq.~\eqref{eq:sed}
converts specific \emph{luminosity} into specific \emph{photon-number} flux
($L_\nu/h\nu$). This kernel is a pure numerical integral over whatever
limits are handed to it — it performs no calibration and needs no
conversion factor, because nothing about it is required to reproduce an
independently-known total.

\subsection{Why its soft/hard limits differ from $C_b$'s}
\label{sec:why-limits}

Both places use different limits for soft vs.\ hard, but the two splits
answer two different physical questions, which is the direct answer to
"why are they needed here but not for \code{integrate\_over\_nu}":

\begin{itemize}
  \item \textbf{$C_b$'s limits} ($[\nu_{\rm th},\nu_{\rm break}]$ vs.\
    $[\nu_{\rm break},\nu_{\rm max}]$, fixed, rest-frame) exist because
    \emph{calibration is band-specific by definition}: $C_b$ is only
    meaningful once "the band" is fixed, since it is derived by demanding
    $\int_{\rm band} L_b(\nu)\,d\nu$ equals that band's own known
    luminosity (Eq.~\eqref{eq:calibration}). Soft and hard are calibrated
    against physically distinct observables (soft-band and hard-band $L_X$,
    with distinct obscuration corrections and, potentially, distinct
    $\alpha_b$), so they necessarily use distinct, fixed limits. This step
    happens once per redshift step, with no $z''$ dependence at all.

  \item \textbf{\code{integrate\_over\_nu}'s limits}
    (\code{lower/upper\_int\_limit\_AGN\_soft/hard}, recomputed every
    $(z',z'')$ pair) exist for an unrelated reason: they track \emph{which
    part of the rest-frame spectrum can physically still be soft or hard by
    the time it arrives at $z'$}. A photon emitted at $z''$ in the rest-frame
    hard band redshifts by $(1+z')/(1+z'')$ on its way to $z'$, and can
    redshift down into what is, at $z'$, the soft band — hence the
    \begin{verbatim}
* (1. + zp) / (1. + zpp)
\end{verbatim}
    factor on every AGN soft/hard limit (\code{ComputeTs.c:793-821}), absent
    from GAL/III's limits because GAL/III use one single band with no
    internal soft/hard split to redshift across. On top of that redshift
    scaling, the lower edge of each band is further clamped to
    $\max(\cdot,\, \nu_{\tau=1}(z',z''))$, the frequency at which the IGM
    becomes optically thick along that shell (\code{nu\_tau\_one\_agn}) —
    a purely radiative-transfer constraint with no analogue in the
    calibration step at all.
\end{itemize}

In short: $C_b$'s limits define \emph{what the band is}, once, in the
rest frame, for calibration purposes. \code{integrate\_over\_nu}'s limits
define \emph{what fraction of each band is currently visible}, recomputed
every shell, for radiative-transfer purposes. They are both "different
limits for soft vs.\ hard" but answer unrelated questions — one is about
the definition of a band, the other about a shifting horizon within it.

\section{Step 4: assembling the heating-rate ODE term}

Per shell $z''_{\rm ct}$ (in \code{evolveInt}, \code{XRayHeatingFunctions.c:1409-1423}):
\begin{align}
  I_b(z'',x_e) &= X_b(z'')\,(1+z'')^{-\alpha_b}\, \mathcal{K}_b(z',z'',x_e), \\
  \frac{dx_{\rm heat}}{dt}\bigg|_b &=
    \mathcal{A}_b(z') \sum_{z''} \frac{dt}{dz''}\,\Delta z''\; I_b(z'',x_e),
  \label{eq:heat-sum}
\end{align}
where $X_b(z'')$ is \code{XAGN\_soft}/\code{\_hard}[z''], the AGN band
emissivity for that shell (an \code{SMOOTHED\_AGN\_(soft/hard)} grid value,
itself built from \code{gal->BHXrayEmissivity} in \code{blackhole\_feedback.c}).
The $(1+z'')^{-\alpha_b}$ factor here is the emission-side counterpart of the
$(1+z')^{\alpha_b+3}$ in $\mathcal{A}_b(z')$ (Eq.~\eqref{eq:prefactor}) — together they implement the same
$(1+z'/1+z'')^{\alpha_b}$ spectral redshifting used to shift the
\code{integrate\_over\_nu} limits in Section~\ref{sec:why-limits}, applied
this time to the amplitude rather than to the integration bounds.
$dt/dz''$ is \code{dtdz(zpp)}, $\Delta z''$ is the shell width
(\code{dzpp}), and the sum runs over \code{TsNumFilterSteps} radial/temporal
shells.

Soft and hard are summed only after each has its own $\mathcal{A}_b$
applied (\code{XRayHeatingFunctions.c:1457-1469}):
\begin{equation}
  \frac{dx_{\rm heat}}{dt}\bigg|_{\rm AGN} =
    \frac{dx_{\rm heat}}{dt}\bigg|_{\rm soft} + \frac{dx_{\rm heat}}{dt}\bigg|_{\rm hard},
\end{equation}
and this combined AGN heating rate enters $dT_K/dz'$ (\code{deriv[3]} /
\code{dansdz[3]}) exactly where the stellar term does, alongside adiabatic
cooling and Compton heating (unchanged from the pre-existing GAL/III
pathway).

\section{Symbol / code cross-reference}

\begin{center}
\begin{tabular}{@{}lll@{}}
\toprule
Symbol & Meaning & Code name \\
\midrule
$\nu_{\rm th}$ & pivot frequency & \code{NuXrayThreshold} \\
$\nu_{\rm break}$ & soft/hard split & \code{NuXraySoftCut} \\
$\nu_{\rm max}$ & hard band ceiling & \code{NuXrayMax} \\
$\alpha_b$ & band spectral index & \code{SpecIndexXrayAGNSoft} / \code{Hard} \\
$L_{{\rm band},b}$ & calibrated band luminosity & \code{quasar\_lx} (obscuration-corrected) \\
$C_b$ & luminosity conversion factor & \code{Luminosity\_converstion\_factor\_AGN\_soft/hard} \\
$\mathcal{A}_b(z')$ & redshift prefactor & \code{const\_zp\_prefactor\_AGN\_soft/hard} \\
$\mathcal{K}_b(z',z'',x_e)$ & frequency shape kernel & \code{freq\_int\_heat\_tbl\_AGN\_soft/hard} \\
$X_b(z'')$ & per-shell AGN emissivity & \code{XAGN\_soft/hard[R\_ct]} \\
$\nu_{\tau=1}(z',z'')$ & opacity horizon & \code{nu\_tau\_one\_agn} \\
\bottomrule
\end{tabular}
\end{center}

\section{Open items for the meeting}

\begin{itemize}
  \item Confirm the sign/exponent convention in Eq.~\eqref{eq:sed}
    ($L_\nu \propto \nu^{-\alpha}$) matches the intended definition of
    \code{SpecIndexXrayAGNSoft/Hard} as used elsewhere in the paper/thesis.
  \item Confirm whether $L_{{\rm band},b}$ (the calibration target) should
    itself already be obscuration-corrected at the point it enters
    Eq.~\eqref{eq:calibration}, or whether obscuration is meant to apply
    only downstream (via \code{NHfrac}/\code{NHTrans}) — this derivation
    assumes the latter (obscuration is applied separately, not baked into
    $C_b$).
  \item Verify the $\alpha_b=1$ branch's tolerance (\code{REL\_TOL}
    $=10^{-5}$, \code{meraxes.h:65}) is tight enough not to introduce a
    discontinuity in $C_b$ near $\alpha_b=1$ for realistic parameter
    sweeps.
\end{itemize}

\end{document}
